% Engendré par saillants.mjs : ne pas modifier à la main.
\section*{Ce que les preuves ajoutent à la page}\addcontentsline{toc}{section}{Ce que les preuves ajoutent à la page}

Les démonstrations qui disent plus que la page, retenues par un critère et non par un choix de l'édition. Une démonstration figure ici si son verdict est \emph{plus général}, \emph{faux sur la page} ou \emph{conjecture de la page réfutée}, et si une piste qu'elle démontre reste ouverte dans la littérature ou a été réfutée : une piste que la recherche a retrouvée imprimée n'est pas comptée. Ce n'est pas un classement par importance, et une piste ouverte n'est pas une affirmation de nouveauté. Les contre-exemples reposent sur les hypothèses telles que la lecture les donne ; celles des cotes où la page est dite fausse restent à relire sur le fac-similé.

{\small\setlength{\tabcolsep}{4pt}
\begin{longtable}{@{}>{\raggedright\arraybackslash}p{1.2cm}>{\raggedright\arraybackslash}p{1.8cm}>{\raggedright\arraybackslash}p{6.2cm}>{\raggedright\arraybackslash}p{2.2cm}p{2.6cm}r@{}}
\toprule
Cote & Pages & Énoncé & Piste & Fichier Lean & \S\\
\midrule
\endhead
\multicolumn{6}{@{}l}{\textit{La page est fausse telle qu'énoncée} (11)}\\[1pt]
160 & 13 & Le déterminant de Moore et les coefficients de Dickson & piste réfutée, piste ouverte & \href{\leanblobmain/Folder160.lean}{\texttt{\detokenize{Folder160}}} & \ref{sec:160}\\
80 & 24 ; 24, 26 & Un contre-exemple au théorème des arrangements cruciaux & piste réfutée & \href{\leanblobmain/Folder80.lean}{\texttt{\detokenize{Folder80}}} & \ref{sec:80}\\
90 & 9 & Un groupe plat irréductible dont un saturé n'est pas de type fini & piste réfutée & \href{\leanblobmain/Folder90.lean}{\texttt{\detokenize{Folder90}}} & \ref{sec:90}\\
63 & 97--99 & Les quatre composantes de $j = 0$ ont quatre tangentes en $p = 2$ & piste réfutée & \href{\leanblobmain/Folder63.lean}{\texttt{\detokenize{Folder63}}} & \ref{sec:63}\\
158 & 21--41, 44, 50--53 & Monoïdes finis, monoïdes à élément minimal, catégories finies & piste ouverte & \href{\leanblobmain/Folder158.lean}{\texttt{\detokenize{Folder158}}} & \ref{sec:158}\\
19 & 4 & Le théorème de comonadicité 1.12, la comonade sur une base produit le théorème de Gabriel--Popescu et les chaînes d'adjoints & piste ouverte & \href{\leanblobmain/Folder19.lean}{\texttt{\detokenize{Folder19}}} & \ref{sec:19}\\
150 & 19 & Voisinages coniques, un contre-exemple sur {[0, 1 & piste ouverte & \href{\leanblobmain/Folder150.lean}{\texttt{\detokenize{Folder150}}} & \ref{sec:150}\\
77 & 85 & Le discriminant universel modulo 2 et modulo 4, la sphère, la symétrie canonique & piste ouverte & \href{\leanblobmain/Folder77.lean}{\texttt{\detokenize{Folder77}}} & \ref{sec:77}\\
29 & 212, 214 & Le critère du déterminant sur l'inertie modérée & piste ouverte & \href{\leanblobmain/Folder29.lean}{\texttt{\detokenize{Folder29}}} & \ref{sec:29}\\
151 & 68 & Le recollement en profondeur deux & piste réfutée & \href{\leanblobmain/Folder151.lean}{\texttt{\detokenize{Folder151}}} & \\
94 & 7 & La réciproque de la remarque 5 bis sous la finitude locale & piste réfutée, piste ouverte & \href{\leanblobmain/Folder94.lean}{\texttt{\detokenize{Folder94}}} & \\
\multicolumn{6}{@{}l}{\textit{Une conjecture de la page est réfutée} (1)}\\[1pt]
156-4 & 65--68 & Les parties saturées ne forment pas une algèbre de Boole & piste réfutée & \href{\leanblobmain/Folder156_4.lean}{\texttt{\detokenize{Folder156_4}}} & \ref{sec:156-4}\\
\multicolumn{6}{@{}l}{\textit{La preuve généralise la page} (21)}\\[1pt]
42 & 3--5 ; 1--2, 5--6 & Familles cohérentes de germes, le lemme 1 et le lemme 2 & piste ouverte & \href{\leanblobmain/Folder42.lean}{\texttt{\detokenize{Folder42}}} & \ref{sec:42}\\
153 & 2--7 & Polynômes à valeurs entières, co-opérations, analyseurs & piste ouverte & \href{\leanblobmain/Folder153.lean}{\texttt{\detokenize{Folder153}}} & \ref{sec:153}\\
161-1 & 2--6 & La partie fixe d'une adjonction, les adjonctions idempotentes, les objets locaux & piste ouverte & \href{\leanblobmain/Folder161_1.lean}{\texttt{\detokenize{Folder161_1}}} & \ref{sec:161-1}\\
39 & 11--14, 18 & Le théorème de préparation « formel » ; fonctions implicites dans une classe de séries & piste ouverte & \href{\leanblobmain/Folder39.lean}{\texttt{\detokenize{Folder39}}} & \ref{sec:39}\\
106 & 9 & Catégories à cofibrations, factorisation de Brown et critère $(\star\star)$ & piste ouverte & \href{\leanblobmain/Folder106.lean}{\texttt{\detokenize{Folder106}}} & \ref{sec:106}\\
67 & 90 & Modules à involution et triples (P, Q, m) & piste ouverte & \href{\leanblobmain/Folder67.lean}{\texttt{\detokenize{Folder67}}} & \ref{sec:67}\\
77 & 68--71, 87, 94--98, 117--124 ; 85 & Le discriminant universel modulo 2 et modulo 4, la sphère, la symétrie canonique & piste ouverte & \href{\leanblobmain/Folder77.lean}{\texttt{\detokenize{Folder77}}} & \ref{sec:77}\\
79 & 19--20, 92--95, 188, 242--247, 275--279 ; 187--188, 275 & Graphes à homothéties, commutant et étoiles sur $(\mathbf Z/n)^2$ & piste ouverte & \href{\leanblobmain/Folder79.lean}{\texttt{\detokenize{Folder79}}} & \ref{sec:79}\\
78 & 109--113 & Le critère de fermeture des polygones réguliers sur un anneau, et les 2p-gones en bande & piste ouverte & \href{\leanblobmain/Folder78.lean}{\texttt{\detokenize{Folder78}}} & \ref{sec:78}\\
40 & 26--27 & La régularité des torseurs sous α\_p ne passe pas à la somme & piste ouverte & \href{\leanblobmain/Folder40.lean}{\texttt{\detokenize{Folder40}}} & \ref{sec:40}\\
136 & 22, 26--31 & Le lemme binomial des ombres & piste ouverte & \href{\leanblobmain/Folder136.lean}{\texttt{\detokenize{Folder136}}} & \ref{sec:136}\\
37 & 92 & Le contre-exemple à EGA I 10.10.5 & piste réfutée, piste ouverte & \href{\leanblobmain/Folder37.lean}{\texttt{\detokenize{Folder37}}} & \ref{sec:37}\\
120 & 8--13, 20--23 & Sommes amalgamées, écraser une partie modérée sur un corps & piste ouverte & \href{\leanblobmain/Folder120.lean}{\texttt{\detokenize{Folder120}}} & \ref{sec:120}\\
145 & 34--35 & L'exposant unité dans le groupe profini libre & piste ouverte & \href{\leanblobmain/Folder145.lean}{\texttt{\detokenize{Folder145}}} & \ref{sec:145}\\
64 & 48 & La configuration octaédrale en caractéristique 5 & piste ouverte & \href{\leanblobmain/Folder64.lean}{\texttt{\detokenize{Folder64}}} & \ref{sec:64}\\
70 & 101--102, 106 & L'icosaèdre en coordonnées universelles, fidèle en toute caractéristique, et sa conique modulo 5 & piste ouverte & \href{\leanblobmain/Folder70.lean}{\texttt{\detokenize{Folder70}}} & \ref{sec:70}\\
128 & 4--9 & Descente le long d'un conducteur nilpotent, idéaux d'équivalence de rang deux & piste ouverte & \href{\leanblobmain/Folder128.lean}{\texttt{\detokenize{Folder128}}} & \ref{sec:128}\\
137 & 77--83 & La pleine fidélité & piste ouverte & \href{\leanblobmain/Folder137.lean}{\texttt{\detokenize{Folder137}}} & \\
5 & 4--6 & L'algèbre artinienne universelle à puissances divisées & piste ouverte & \href{\leanblobmain/Folder5.lean}{\texttt{\detokenize{Folder5}}} & \ref{sec:5}\\
148 & 55--58 & Unicité du recollement & piste ouverte & \href{\leanblobmain/Folder148.lean}{\texttt{\detokenize{Folder148}}} & \\
154 & 153 & Cube des générisations, sous la forme précisée par l'édition & piste ouverte & \href{\leanblobmain/Folder154.lean}{\texttt{\detokenize{Folder154}}} & \\
\bottomrule
\end{longtable}}
